
\section{Schéma Monolithique (DG/FV)}
La méthode se repose sur deux détails clés : une réécriture du schéma DG utilisant des sous-mailles et 
un flux hybride entre le flux DG du schéma réécrit et un flux Volume Fini entre sous-mailles.
\subsection{Réecriture du schéma DG}
\paragraph{sous-maillage \\}

On dote chaque maille $\omega_i$ d'un sous-maillage : $\{S_m^i\}_m = ]\tilde{x}^i_{m-1/2}, \tilde{x}^i_{m+1/2}[$ vérifiant les mêmes propriétés que le maillage de $\mathbb{D}$ : \eqref{maillage1}-\eqref{maillage3}\\
On cherchera à passer les degrées de libertés des polynomes de $ \sigma \in\mathbb{P}^k$, des moments sur la base polynomiale $ \left\{ \underline{\sigma}_p \right\}_p$ aux moyennes de ces polynomes sur le sous-mailles $\{S_m^i \}$ : $ \left\{\overline{\sigma}^m \right\}_m$ \\
\noindent Soit la matrice $$\mathcal{P}_{m,p} = \frac{1}{|S_m^i|} \int_{S_m^i} \sigma_p(x) dx$$ 
On remarque que pour que le passage des moments polynomiaux $\left\{ \underline{\sigma}_p \right\}_p$ aux sous-volumes finis soit injectif, il nous faut : $N_s \geq N_k$. 
De manière générale, on dit que le sous-maillage est admissible si et seulement si la matrice $\mathcal{P}^t\,\mathcal{P}$ est inversible.

\paragraph{Remarque \\}
On note que pour avoir une équivalence au sens le plus strict possible, il nous faut $N_s = N_k$. 
Dans le cas échéant ($N_s \geq N_k$), on utilisera un opérateur de reconstruction $\mathcal{R} = \left(\mathcal{P}^t\,\mathcal{P}\right)^{-1} \mathcal{P}^t$  qui nous donnera les moments polynomiaux au sens des moindres carrés.
\newline

Il est alors simple de passer d'une représentation à une autre : 
\begin{equation}
    \overline{U} = \mathcal{P}\underline{U} \qquad \underline{U} = \mathcal{R}\overline{U}
\end{equation}

On en profite pour introduire les fonctions $\{\phi_m^i\}_{1\leq m\leq N_s}$ les projections $\mathbb{P}^k$ de $\mathbb{1}_{S^i_m}$ : 
$\int_{\omega_i} \phi_m \sigma dx = \int_{S_m^i} \sigma dx$, $\forall \sigma \in \mathbb{P}^k(\omega_i)$

\paragraph{ Construction des flux reconstruit \\}

En reprenant la formulation variationnelle \eqref{Formulation_Variationnelle}, si l'on projette $\mathbf{f}(\mathbf{u_h})$ sur $\mathbb{P}^k$ : $\mathbf{f_h}$, on obtient une formulation équivalente : 
$$\int_{\omega_i} \partial_t \mathbf{u_h} \sigma dx - \int_{\omega_i} \mathbf{f}_h^i \partial_x \sigma dx 
+ \left[ \mathcal{F}\, \sigma \right]_{x_{i-\frac{1}{2}}}^{x_{i+\frac{1}{2}}} = 0 \qquad \forall  \sigma \in \mathbb{P}^k(\omega_i) $$  
Alors par le biais d'une intégration par partie, on obtient la forme forte des schéma DG :  
$$\int_{\omega_i} \partial_t \mathbf{u_h} \sigma dx = \int_{\omega_i} \partial_x \mathbf{f}_h^i  \sigma dx 
- \left[ ( \mathbf{f}_h^i - \mathcal{F} )\, \sigma \right]_{x_{i-\frac{1}{2}}}^{x_{i+\frac{1}{2}}} \qquad \forall  \sigma \in \mathbb{P}^k(\omega_i) $$  

Comme cette équation est vrai pour tout polynôme de $\mathbb{P}^k(\omega_i)$, il est en particulier vrai pour $\phi_m$. Cela nous amène presque à un flux ne dépendant que de $S_m^i$ : 
\begin{equation}
    \int_{S_m^i} \partial_t \mathbf{u_h} dx = - \left(\int_{S_m^i}  \partial_x \mathbf{f}_h^i dx - \left[ ( \mathbf{f}_h^i - \mathcal{F} )\, \phi_m \right]_{x_{i-\frac{1}{2}}}^{x_{i+\frac{1}{2}}}\right) \qquad \forall m\in \{ 1, \dots, N_s\}
\end{equation}
Il est cependant possible d'obtenir un schéma volume fini sur les $\{\overline{u_h}^m\}$ en fabriquant un flux avec pour unique but de ressembler à un flux volume fini : \\
Soit $\widehat{F^i}_{m+1/2}$ ces flux reconstruits posé telle que : 

\begin{align}
&\widehat{F^i}_{m+1/2} - \widehat{F^i}_{m-1/2} =  -  \left[\mathbf{f}_h^i\right]_{\tilde{x}^i_{m-1/2}}^{\tilde{x}^i_{m+1/2}}  - \left[ ( \mathbf{f}_h^i - \mathcal{F} )\, \phi_m \right]_{x_{i-\frac{1}{2}}}^{x_{i+\frac{1}{2}}} \qquad \forall m\in \{ 1, \dots, N_s\} \\
& \text{où l'on impose : } \widehat{F^i}_{1/2} = \mathcal{F}_{i-1/2} \qquad \widehat{F^i}_{ N_s + 1/2} = \mathcal{F}_{i+1/2}
\end{align}
Il est possible d'obtenir une forme close de $\widehat{F^i}_{m+1/2}$: 
\begin{align}
    &\widehat{F^i}_{m+1/2} = \mathbf{f}_h^i(\tilde{x}^i_{m+1/2}) 
    -C^{i-1/2}_{m+1/2} ( \mathbf{f}_h^i(x_{i-1/2}) - \mathcal{F}_{i-1/2}) 
    -C^{i+1/2}_{m+1/2} ( \mathbf{f}_h^i(x_{i+1/2}) - \mathcal{F}_{i+1/2}) \\
    & \text{où }  C^{i-1/2}_{m+1/2} = (1- \sum_{p=1}^{m}\phi_p(x_{i-1/2})) \, , \qquad
        C^{i+1/2}_{m+1/2} =     \sum_{p=1}^{m}\phi_p(x_{i+1/2})
\end{align}

Tout ceci nous permet d'écrire le schéma DG sous une forme Volumes Finis : 
 
\begin{equation}\label{DG_as_FV}
    \partial_t \overline{\mathbf{u}_m}^i = -\frac{1}{|S_m^i|}(\widehat{F^i}_{m+1/2} - \widehat{F^i}_{m-1/2})
\end{equation}

\paragraph{Remarque \\}
\noindent Même si la forme de \eqref{DG_as_FV} laisse à penser que l'évolution de $\overline{\mathbf{u_h}}^m$ 
ne dépend uniquement des sous-cellules voisines à $S_m^i$, 
il faut se rappeler que $\widehat{F^i}_{m+1/2}$ dépend non seulement des sous-cellules voisines, 
mais aussi des cellules voisines $\omega_{i\pm1}$ par le biais du flux numérique évalué aux interfaces des cellules $x_{i\pm1/2}$ \\

Les flux $\widehat{F^i}_{m+1/2}$ pouvant se trouver entre deux sous-cellules appartenant à la même maille ou de mailles différentes, 
Par soucis de claireté, nous écrirons $\widehat{F_{mp}}$ le flux allant de la sous-cellule $S_m^i$ à la sous-cellule $S_p^j$, $j\in \{i-1,i,i+1\}$ 

\subsection{Flux hybride (DG/FV)}

Maintenant que nous établie un schéma agissant sur les sous valeurs moyennes, 
équivalent au schéma DG agissant sur les moments polynomiaux. 
Il devient alors possible de combiner les flux DG $\widehat{F_{mp}}$ 
avec des flux proprement volumes finis défini sur les sous-mailles $F^{\text{FV}}_{mp}$. 

\begin{equation} 
\begin{aligned}
    \widetilde{F_{mp}}  &= F^{\text{FV}}_{mp} + \theta_{mp}\,\, (\widehat{F_{mp}} - F^{\text{FV}}_{mp})\\
                        &= F^{\text{FV}}_{mp} + \theta_{mp}\,\, \Delta F_{mp}
\end{aligned}
\end{equation}

Le schéma dépend alors du choix de $\theta_{mp}$. 
Si l'on pose $\theta_{mp} = 1$, alors nous obtiendrons un schéma DG qui pourra être d'ordre élevée.
Si l'on pose $\theta_{mp} = 0$, alors nous obtiendrons un schéma Volumes Finis qui sera entropique et suivra le principe du maximum local. \\
Avant de s'attarder sur le choix des $\theta_mp$, regardons un minimum le schéma que nous avons obtenus : 

\begin{equation}\label{DG_as_FV}
    \partial_t \overline{\mathbf{u}_m}^i =  - \frac{1}{|S_m^i|} \sum_{S_p^j \in \mathcal{V}_m^i} \widetilde{F_{mp}}
\end{equation}

En se souvenant que toutes propriétés d'un schéma \textit{RK-SSP} d'ordre élevée est dérivée des propriétés du schéma d'Euler explicite, 
il est possible de vois que $\overline{\mathbf{u}_m}^{i,n+1}$ peut se voir comme une combinaison convexe entre $\overline{\mathbf{u}_m}^{i,n}$ et ce que l'on définira comme des états de Riemann intermédiaires hybrides : 
$ \widetilde{\mathbf{u}}_{mp}^- = \overline{\mathbf{u}_m}^{i,n} - \frac{\widetilde{F_{mp}} -  F_h^i(\overline{\mathbf{u}_m}^{i,n})\cdot n_{mp} }{\gamma_{mp}}$

\begin{align*}
&\overline{\mathbf{u}_m}^{i,n+1} = \overline{\mathbf{u}_m}^{i,n} - \frac{\Delta t_n}{|S_m^i|} \sum_{S_p^j \in \mathcal{V}_m^i} \widetilde{F_{mp}} 
\pm \frac{\Delta t_n}{|S_m^i|} \sum_{S_p^j \in \mathcal{V}_m^i} \gamma_{mp}\overline{\mathbf{u}_m}^{i,n} \pm \frac{\Delta t_n}{|S_m^i|} F_h^i(\overline{\mathbf{u}_m}^{i,n}) \\
&\overline{\mathbf{u}_m}^{i,n+1}= \left( 1-  \frac{\Delta t_n}{|S_m^i|} \sum_{S_p^j \in \mathcal{V}_m^i} \gamma_{mp}\right) \overline{\mathbf{u}_m}^{i,n} 
                                + \frac{\Delta t_n}{|S_m^i|} \sum_{S_p^j \in \mathcal{V}_m^i} \gamma_{mp} \left( \overline{\mathbf{u}_m}^{i,n} - \frac{\widetilde{F_{mp}} -  F_h^i(\overline{\mathbf{u}_m}^{i,n})\cdot n_{mp} }{\gamma_{mp}} \right)
\end{align*}
On note ici que la combinaison est convexe seulement si $ \frac{\Delta t_n}{|S_m^i|} \sum_{S_p^j \in \mathcal{V}_m^i} \gamma_{mp} \leq 1$, ce qui nous donne une condition CFL : 

\begin{equation}
    \Delta t \leq \frac{|S_m^i|}{\sum_{S_p^j \in \mathcal{V}_m^i} \gamma_{mp}}
\end{equation} 
Elle sera par ailleurs la seul condition CFL à prendre en compte car plus stricte que la condition CFL du schéma DG.\\ 

Il sera souvent plus simple d'étudier les combinaisons convexes de $\overline{\mathbf{u}_m}^{i,n+1}$ en découpant $ \widetilde{\mathbf{u}}_{mp}^-$ comme : 
\begin{equation}
\begin{aligned}
    \widetilde{\mathbf{u}}_{mp}^- =& \left( \frac{\overline{\mathbf{u}_m}^{i,n}+ \overline{\mathbf{u}_p}^{j,n}}{2} - \frac{(F_h^i(\overline{\mathbf{u}_p}^{j,n}) - F_h^i(\overline{\mathbf{u}_m}^{i,n})   )\cdot n_{mp}}{2 \gamma_{mp}}   \right) - \theta_{mp} \frac{\Delta F_{mp}}{\gamma_{mp}} \\
    =& \mathbf{u}_{mp}^{*,FV} - \theta_{mp} \frac{\Delta F_{mp}}{\gamma_{mp}}
\end{aligned}    
\end{equation}

On note au passage ces égalité de symmétrie/ d'antisymmétrie : \\
$n_{mp}= -n_{pm}$, et donc : $F^{\text{FV}}_{mp} = - F^{\text{FV}}_{pm}$, $\widehat{F_{mp}} = -\widehat{F_{pm}}$ et $\widetilde{F_{mp}} = -\widetilde{F_{pm}}$\\
tandis que l'on pose $\theta_{pm} = \theta_{mp}$, ce qui nous donne : $\widetilde{\mathbf{u}}_{pm}^- = \widetilde{\mathbf{u}}_{mp}^+ = \mathbf{u}_{mp}^{*,FV} + \theta_{mp} \frac{\Delta F_{mp}}{\gamma_{mp}}$, car $ \mathbf{u}_{mp}^{*,FV} = \mathbf{u}_{pm}^{*,FV}$ \\

Tout ceci nous ammène donc à la remarque suivante : si l'on souhaite que la solution se trouve dans un ensemble convexe $\Omega$, tant que $\overline{\mathbf{U}}^{n} \in \Omega$, 
il est toujours possible de trouver $\theta_{mp}$ pour que $\overline{\mathbf{U}}^{n+1} \in \Omega$ car $\mathbf{u}_{mp}^{*,FV}$ est combinaison convexe de $ \overline{\mathbf{u}}_{mp}^{n}$ {\color{red} Voir annexe}.

\subsection{choix du paramètre d'hybridation {\color{red} nom à changer}}

\paragraph{Global Maximal Principle (GMP)}
La première propriété que nous allons étudier et qui nous servira d'exemple pour la suite est comment imposer un principe du maximum global dans le cas scalaire $\Omega \subset \R$, $M=1$ : 
$$ \alpha \leq \overline{U}^{0} \leq \beta \implies \alpha \leq \overline{U}^{n} \leq \beta \quad \forall n$$

\paragraph{Local  Maximal Principle (LMP)}
\paragraph{positivité d'un schéma sur le problème de la dynamique des gaz}
\paragraph{Détections d'extrema locaux}